\documentclass[10pt,a4paper]{amsart}
\usepackage[margin=19mm]{geometry}
\usepackage{amsmath,amssymb,mathtools}
\usepackage[scr=rsfs,cal=euler]{mathalfa}
\usepackage{xfrac}
\usepackage[unicode,hidelinks,bookmarks=false]{hyperref}
\newcommand{\ie}{i.e.\ }
\newcommand{\cf}{cf.\ }
\newcommand{\resp}{respectively\ }
\newcommand{\Subsection}{\subsection}
\providecommand{\nobreakdash}{\nobreak\hbox{-}}
\setlength{\emergencystretch}{2em}
\multlinegap0pt
\usepackage{amssymb}
\usepackage{mathtools}
\usepackage[shortlabels]{enumitem}
\newtheorem{proposition}{Proposition}
\newcommand{\E}{\mathcal{E}}
\title{Lifted Takahashi Convexity on the Isbell-Convex Hull of an Asymmetrically Normed Real Vector Space}
\author{Philani Rodney Majozi, Mcedisi Sphiwe Zweni}
\date{Source: arXiv:2603.12101v3 (CC BY 4.0)}
\begin{document}
\maketitle

\noindent\emph{Source and licence.} Lifted Takahashi Convexity on the Isbell-Convex Hull of an Asymmetrically Normed Real Vector Space. Version arXiv:2603.12101v3 (CC BY 4.0), distributed by arXiv under the Creative Commons Attribution 4.0 International licence, http://creativecommons.org/licenses/by/4.0/, as recorded in the arXiv licence metadata for this exact version and shown on the abstract page https://arxiv.org/abs/2603.12101v3. Full licence notice: \texttt{licenses/arxiv-2603.12101-CC-BY-4.0.txt}.\par\smallskip

\noindent\textit{Editorial scope.} Counted unit: the article's proposition prop:oplusWconvex (source0.tex lines 1299--1304). The article's complete original proof of it is delivered with it (source0.tex lines 1306--1395, one \texttt{proof} environment of 90 lines). No mathematics is written, repaired or re-derived: the statement and the proof are the article's own lines, in the article's own notation, and no retained line is substituted or re-flowed.  No further source-local context is needed: every cross-reference of every retained line resolves inside the two delivered spans.  Nothing was omitted from any retained span, so the package carries no omission record.  No retained line contains a citation, so the wrapper carries no bibliography and no bibliography entry.  No retained line contains a figure environment, an image-inclusion command or drawing code, so no image is delivered and the package declares no asset dependency.  Editorial formatting only, in addition: an AMS \texttt{amsart} preamble that restates the article's own declarations and macros used by the retained lines, namely the environments \texttt{proposition} on separate counters, together with the standard packages those lines need.  The retained environments print in this document as Proposition 1 in order of appearance, whereas the article prints the same items with the numbering of its own full document.  The complete native source of the article as distributed in the arXiv e-print for this exact version is bundled unchanged as \texttt{sources/196259/source0.tex}.
\begin{proposition}\label{prop:oplusWconvex}
Let \((X,\|\cdot\|,W)\) be a convex asymmetrically normed real vector space.
Assume that \(W\) is translation-invariant.  If
\(f,g\in\E(X,\|\cdot\|)\) are \(W\)-convex on \(X\), then
\(f\oplus g\) is \(W\)-convex on \(X\).
\end{proposition}

\begin{proof}
Let \(f=(f_1,f_2)\) and \(g=(g_1,g_2)\).  For \(j=1,2\), set
\[
j' = 3-j.
\]
Thus \(1'=2\) and \(2'=1\).  The hull addition may be written uniformly as
\[
(f\oplus g)_j(x)
=
\sup_{z\in X}
\bigl(f_j(x-z)\dot{-}g_{j'}(z)\bigr).
\]
Indeed, for \(j=1\) this says
\[
(f\oplus g)_1(x)
=
\sup_{z\in X}
\bigl(f_1(x-z)\dot{-}g_2(z)\bigr),
\]
while for \(j=2\) it says
\[
(f\oplus g)_2(x)
=
\sup_{z\in X}
\bigl(f_2(x-z)\dot{-}g_1(z)\bigr).
\]

Fix \(x,y\in X\) and \(\lambda\in[0,1]\).  Translation invariance gives
\[
W(x,y,\lambda)-z
=
W(x-z,y-z,\lambda)
\qquad (z\in X).
\]
Therefore
\[
(f\oplus g)_j\bigl(W(x,y,\lambda)\bigr)
=
\sup_{z\in X}
\bigl(
f_j(W(x,y,\lambda)-z)\dot{-}g_{j'}(z)
\bigr)
\]
\[
=
\sup_{z\in X}
\bigl(
f_j(W(x-z,y-z,\lambda))\dot{-}g_{j'}(z)
\bigr).
\]
Since \(f_j\) is \(W\)-convex,
\[
f_j(W(x-z,y-z,\lambda))
\le
\lambda f_j(x-z)+(1-\lambda)f_j(y-z).
\]
Hence
\[
(f\oplus g)_j\bigl(W(x,y,\lambda)\bigr)
\le
\sup_{z\in X}
\Bigl(
\lambda f_j(x-z)+(1-\lambda)f_j(y-z)\dot{-}g_{j'}(z)
\Bigr).
\]

Using the elementary inequality
\[
\bigl(\lambda A+(1-\lambda)B-C\bigr)^+
\le
\lambda(A-C)^+ + (1-\lambda)(B-C)^+
\]
for \(A,B,C\in\mathbb R\), we obtain
\[
(f\oplus g)_j\bigl(W(x,y,\lambda)\bigr)
\le
\lambda
\sup_{z\in X}\bigl(f_j(x-z)\dot{-}g_{j'}(z)\bigr)
+
(1-\lambda)
\sup_{z\in X}\bigl(f_j(y-z)\dot{-}g_{j'}(z)\bigr).
\]
Therefore
\[
(f\oplus g)_j\bigl(W(x,y,\lambda)\bigr)
\le
\lambda (f\oplus g)_j(x)+(1-\lambda)(f\oplus g)_j(y).
\]
Since this holds for \(j=1,2\), the pair \(f\oplus g\) is \(W\)-convex.
\end{proof}

\end{document}
